\subsubsection{Montaje y condiciones del ensayo}

Para contrastar el modelo térmico de la Sección~\ref{sec:modelo_dinamico} se realizó un ensayo en lazo abierto sobre el tanque de mezcla del prototipo, con agua y sin controlador, aplicando escalones de potencia sobre la resistencia de inmersión E-101. El volumen de agua del ensayo se estimó en aproximadamente 18\,L, ya que no pudo medirse de forma directa; este valor se adoptó únicamente para el análisis de los datos del ensayo. La temperatura se midió con el transmisor TT-101 (sensor PT100 con transmisor de 4--20\,mA). La potencia al 100\,\% del ciclo de trabajo se obtuvo midiendo la tensión (126,3\,V) y la corriente (16,2\,A) con multímetro y pinza amperimétrica, lo que dio una potencia medida de $P_{\mathrm{on}} = 2046\,\mathrm{W}$, frente al valor nominal de 2000\,W empleado en el resto del documento.

La potencia se moduló con control proporcional en el tiempo mediante el SSR, programado en una interrupción cíclica (OB30) de 100\,ms. El periodo de modulación fue de 10\,s (100 ciclos), con una resolución de 1\,\% en el ciclo de trabajo. Para conocer la potencia realmente aplicada, el programa acumula el tiempo en que el SSR permanece activo ($t_{\mathrm{on}}$), de modo que la potencia media en cualquier intervalo se calcula, sin depender del valor comandado, como:

\begin{equation}
P = P_{\mathrm{on}}\,\frac{\Delta t_{\mathrm{on}}}{\Delta t}
\label{eq:potencia_media}
\end{equation}

Como protección propia del ensayo, el programa incluye un aborto enclavado por software ante una temperatura superior a 65\,°C o ante la apertura del lazo de 4--20\,mA.

\subsubsection{Adquisición de datos}

El registro se realizó con la función DataLog del PLC S7-1200, que guarda los datos en un archivo CSV dentro de la memoria de carga de la CPU, sin necesidad de un computador conectado durante la prueba. Esta función emplea cuatro instrucciones asíncronas que se ejecutan en secuencia:

\begin{itemize}
\item \texttt{DataLogCreate}: crea el archivo con su encabezado y el número máximo de registros.
\item \texttt{DataLogOpen}: abre el archivo.
\item \texttt{DataLogWrite}: añade un registro con marca de tiempo cada vez que su entrada \texttt{REQ} recibe un flanco ascendente.
\item \texttt{DataLogClose}: cierra el archivo.
\end{itemize}

El archivo se descarga después desde la página \textit{Data logs} del servidor web integrado del PLC. La secuencia se implementó en el OB1 como una máquina de estados (reposo, creación, apertura, listo, escritura, cierre y error). La captura de cada muestra se realiza en el OB30 cada 5\,s (50 ciclos): se copian en la estructura de datos los valores del mismo instante y se activa una bandera que dispara la escritura, la cual se desactiva al terminar. Cada registro contiene el reloj del ensayo (ms), la temperatura (°C), el ciclo de trabajo aplicado (\%) y el tiempo acumulado en ON (ms); en la versión actual del programa se agrega además la consigna de temperatura y el archivo admite un máximo de 1000 registros (valores preliminares, Sección~\ref{sec:montaje_programacion}).

\subsubsection{Procedimiento experimental}

Con el agua en equilibrio con el ambiente ($T_{\mathrm{amb}} \approx 23{,}5\,°\mathrm{C}$, tomada como el promedio de la línea base previa al calentamiento), se ejecutó la secuencia de la Tabla~\ref{tab:escalones_potencia}, cuyos tiempos son aproximados sobre el eje del registro. La variación de la potencia entre escalones proporciona la excitación necesaria para identificar la capacidad térmica y las pérdidas con un ensayo mucho más corto que la constante de tiempo del proceso \cite{seborg2016process}. La temperatura máxima alcanzada fue cercana a 60\,°C, inferior al límite de aborto.

\begin{table}[H]
\centering
\caption{Secuencia de escalones de potencia aplicados durante el ensayo.}
\label{tab:escalones_potencia}
\begin{tabular}{clcc}
\toprule
Segmento & Etapa & Intervalo (min) & Ciclo de trabajo (\%) \\
\midrule
0 & Línea base         & 0 -- 2   & 0   \\
1 & Calentamiento 1    & 2 -- 11  & 100 \\
2 & Calentamiento 2    & 11 -- 24 & 50  \\
3 & Calentamiento 3    & 24 -- 34 & 100 \\
4 & Enfriamiento libre & 34 -- 54 & 0   \\
\bottomrule
\end{tabular}
\end{table}

\subsubsection{Estimación de parámetros}

El modelo de primer orden que describe la dinámica térmica del sistema,

\begin{equation}
C\,\frac{dT}{dt} = P(t) - UA \left[T(t) - T_{\mathrm{amb}}\right],
\label{eq:balance_energia}
\end{equation}

se identificó en MATLAB con tres métodos, todos basados en la potencia real obtenida a partir de $t_{\mathrm{on}}$ (Ecuación~\ref{eq:potencia_media}):

\begin{itemize}
\item \textbf{Método A:} la Ecuación~\ref{eq:balance_energia} se reescribe en la forma lineal en los parámetros
\begin{equation}
\frac{dT}{dt} = a\,P(t) - b\,\Delta T(t), \qquad a = \frac{1}{C}, \qquad b = \frac{UA}{C},
\label{eq:regresion_lineal}
\end{equation}
donde $\Delta T(t) = T(t) - T_{\mathrm{amb}}$, y se resuelve por mínimos cuadrados sobre ventanas deslizantes de 60\,s (múltiplo del periodo de modulación) con los datos del calentamiento.
\item \textbf{Método B:} se minimiza el error entre la temperatura medida y la simulada con la solución exacta del modelo, usando solo los datos del calentamiento.
\item \textbf{Método C:} igual que el método B, pero con todo el registro, excluyendo los primeros 120\,s del enfriamiento por el transitorio térmico de la resistencia.
\end{itemize}

Los métodos A y B se validaron con el tramo de enfriamiento, que no se usó en su ajuste. A partir de $C$ y $UA$, la constante de tiempo y la ganancia del modelo se obtienen como:

\begin{equation}
\tau = \frac{C}{UA}, \qquad K = \frac{1}{UA}.
\label{eq:parametros_recuperados}
\end{equation}

\subsubsection{Parámetros obtenidos}

La Tabla~\ref{tab:resultados_metodos} presenta los parámetros obtenidos con cada método. Se adoptó el método C por usar todo el registro y tener el intervalo de confianza más estrecho.

\begin{table}[H]
\centering
\caption{Parámetros del modelo térmico obtenidos con cada método.}
\label{tab:resultados_metodos}
\begin{tabular}{lccccc}
\toprule
Método & $C$ (J/K) & $UA$ (W/K) & $\tau$ (min) & $K$ (°C/W) & RMSE de validación (°C) \\
\midrule
A & 71\,390 & 14,18 & 83,9 & 0,0705 & 0,060 \\
B & 73\,442 & 12,26 & 99,8 & 0,0816 & 0,573 \\
C & 72\,330 & 13,63 & 88,4 & 0,0733 & 0,163 \\
\bottomrule
\end{tabular}
\end{table}

La Figura~\ref{fig:control_ident_temperatura} compara la temperatura medida con la simulada por los tres modelos durante todo el ensayo, y la Figura~\ref{fig:control_ident_validacion}, la predicción de cada modelo en el tramo de enfriamiento, que no se usó en el ajuste de los métodos A y B.

\insertarfigura[0.85\textwidth]{control_ident_temperatura.pdf}{Temperatura medida durante el ensayo y temperatura simulada con los modelos identificados por los métodos A, B y C.}{fig:control_ident_temperatura}

\insertarfigura[0.85\textwidth]{control_ident_validacion.pdf}{Validación de los modelos con el tramo de enfriamiento libre (potencia nula).}{fig:control_ident_validacion}

La Tabla~\ref{tab:comparacion_parametros} compara los parámetros experimentales con los teóricos del modelo del primer corte, calculados con el volumen de agua del ensayo (18\,L).

\begin{table}[H]
\centering
\caption{Comparación entre parámetros teóricos y experimentales del modelo térmico.}
\label{tab:comparacion_parametros}
\begin{tabular}{lccc}
\toprule
Parámetro & Valor teórico (18\,L) & Valor experimental (método C) & Diferencia (\%) \\
\midrule
$UA$ (W/K)    & 4,00       & 13,63      & $+240{,}9$ \\
$C$ (J/K)     & 75\,348    & 72\,330    & $-4{,}0$ \\
$\tau$ (min)  & 313,9      & 88,4       & $-71{,}8$ \\
$K$ (°C/W)    & 0,250      & 0,0733     & $-70{,}7$ \\
\bottomrule
\end{tabular}
\end{table}

Se estableció como criterio de aceptación una diferencia relativa máxima del $15\,\%$ entre el valor teórico y el experimental de $UA$:

\begin{equation}
\left| \frac{UA_{\mathrm{exp}} - UA_{\mathrm{teo}}}{UA_{\mathrm{teo}}} \right| \times 100\,\% \leq 15\,\%.
\label{eq:criterio_aceptacion}
\end{equation}

Como la diferencia en $UA$ supera este criterio, el modelo teórico no queda confirmado y se actualiza con los valores experimentales: las pérdidas térmicas reales son aproximadamente 3,4 veces las supuestas inicialmente. El modelo adoptado, con la potencia como entrada, es:

\begin{equation}
G(s) = \frac{0{,}0733}{5305\,s + 1}\quad [°\mathrm{C/W}],
\label{eq:modelo_experimental}
\end{equation}

y, con el ciclo de trabajo del SSR como entrada:

\begin{equation}
G_p(s) = \frac{1{,}50}{5305\,s + 1}\quad [°\mathrm{C}/\%].
\label{eq:modelo_experimental_duty}
\end{equation}

El modelo también se extrapoló al volumen de diseño de 20\,L, con lo que $\tau$ pasaría a 5919\,s (98,7\,min). Esa extrapolación no se adopta, porque la capacidad térmica identificada es menor que la del agua sola, lo que impide separar la masa térmica del recipiente y escalar el modelo con confianza. Aun así, los 5919\,s quedan dentro de la envolvente de $\tau$ que se usa en el diseño del controlador (Sección~\ref{sec:diseno_pi}).

\subsubsection{Discusión y limitaciones}

\begin{itemize}
\item El volumen de agua se estimó, no se midió. La capacidad térmica identificada resulta un 4\,\% inferior a la del agua sola, diferencia que cae dentro de esa incertidumbre.
\item El modelo de primer orden no incluye retardo. En los datos se observa un retardo aparente de unos 30\,s al encender: la temperatura permanece dentro del ruido ($\pm 0{,}06\,^{\circ}\mathrm{C}$) durante los primeros 30\,s y a los 35\,s ha subido 0,40\,°C, cuando el modelo predice cerca de 1\,°C. Se atribuye al transporte de calor en agua sin agitación y a la inercia del sensor. Sin embargo, al estimar el retardo por mínimos cuadrados sobre todo el registro, el valor óptimo es de 0 a 5\,s (RMSE de 0,19\,°C sin retardo, frente a 0,59\,°C con 35\,s), porque el ajuste global compensa ese retraso inicial con los demás parámetros. Por ello, en el diseño del controlador se adopta un retardo de 35\,s como cota superior conservadora, y no como valor identificado. Este valor es menor al 1\,\% de $\tau$ (0,66\,\%).
\item Las pérdidas no son del todo lineales, ya que la convección y la evaporación crecen con la temperatura; por ello, $UA$ es un valor efectivo para el rango de 23 a 60\,°C. Esto explica la peor validación del método B.
\item Los tres métodos difieren cerca de un 14\,\% en $UA$, y esa dispersión se toma como incertidumbre del modelo. La sensibilidad a $T_{\mathrm{amb}}$ es baja: una variación de $\pm 2\,°\mathrm{C}$ cambia $UA$ en cerca de $\pm 4\,\%$.
\end{itemize}
